\documentclass[11pt]{article}
\usepackage[margin=0.9in]{geometry}
\usepackage[T1]{fontenc}
\usepackage{lmodern,amsmath,amssymb,amsthm,microtype,booktabs,needspace}
\usepackage[colorlinks=true,linkcolor=blue,urlcolor=blue,citecolor=blue]{hyperref}
\newtheorem{theorem}{Theorem}
\newtheorem{lemma}[theorem]{Lemma}
\newcommand{\N}{\mathbb N}
\newcommand{\ind}{\mathbf 1}
\title{Exact long equality runs in a greedy sequence\protect\\avoiding four-term sums}
\author{Mark Vladimir Prizmic}
\date{November 27, 2026\quad Version 0.1}
\begin{document}
\maketitle
\begin{abstract}
Let $(a_n)$ be the increasing greedy sequence starting at $1$ that excludes
sums of four consecutive previously selected terms. We prove that its maximal
runs of equality with $\lfloor(5n-1)/4\rfloor$, of length at least four, are
exactly the intervals
\[
 \left[\frac{7\cdot16^m-2}{5},\frac{12\cdot16^m+8}{5}\right],
 \qquad m\ge0.
\]
Their lengths are $16^m+3$. This establishes the length formula recorded as
Conjecture~2 in OEIS A394676, with corrected endpoints for several displayed
examples. The proof gives an exact recursion for the extreme values of the
complementary sequence of forbidden sums. All arguments needed for the result
are included.
\end{abstract}

\section{The sequence and the result}

Set $a_1=1$. After $a_1,\ldots,a_n$ have been selected, choose $a_{n+1}$ as
the smallest integer greater than $a_n$ outside
\[
 \{a_j+a_{j+1}+a_{j+2}+a_{j+3}:1\le j\le n-3\}.
\]
The forbidden set is empty for $n<4$. Consecutive refers to positions in the
selected sequence. This is Hoang Xuan Thanh's sequence A394676~\cite{oeis}.
Put
\[
 U_n=\left\lfloor\frac{5n-1}{4}\right\rfloor,
 \qquad b_j=\sum_{t=0}^3a_{j+t}.
\]
An \emph{equality run} is an interval of consecutive indices on which $a_n=U_n$;
it is maximal if it cannot be extended in either direction.

\begin{theorem}\label{thm:runs}
The maximal equality runs of length at least four are exactly
\begin{equation}\label{eq:runs}
 [S_m,E_m]_{\mathbb Z},\qquad
 S_m=\frac{7\cdot16^m-2}{5},\qquad
 E_m=\frac{12\cdot16^m+8}{5},\qquad m\ge0.
\end{equation}
Their lengths are $E_m-S_m+1=16^m+3$. Every other maximal equality run has
length at most three.
\end{theorem}

The first five runs are as follows.
\begin{center}
\begin{tabular}{rrrr}
\toprule
$m$ & First index & Last index & Length\\
\midrule
0&1&4&4\\
1&22&40&19\\
2&358&616&259\\
3&5\,734&9\,832&4\,099\\
4&91\,750&157\,288&65\,539\\
\bottomrule
\end{tabular}
\end{center}
The retrieved OEIS entry gives the same length formula, but its third, fourth,
and fifth displayed intervals have both endpoints one too large. The
qualification ``length at least four'' is necessary: shorter maximal runs
also occur.

\section{Complementarity and uniform bounds}

The next lemma is the four-term case of the companion general-bound
note~\cite{bounds}. We include its proof so the present argument is complete.

\begin{lemma}\label{lem:bounds}
The increasing sequences $(a_n)$ and $(b_j)$ partition the positive integers.
For every $j\ge1$,
\begin{equation}\label{eq:b-bound}
 5j\le b_j\le5j+5.
\end{equation}
\end{lemma}

\begin{proof}
Only finitely many integers are forbidden at any finite stage, so the
construction continues indefinitely. Every $b_j$ exceeds its largest summand
$a_{j+3}$, and becomes forbidden when that summand is selected, before the
greedy construction can reach $b_j$. Conversely, an omitted positive integer
must already have been forbidden when the greedy construction skipped it.
Also $b_{j+1}-b_j=a_{j+4}-a_j\ge4$, so the forbidden sums are distinct. This
proves complementarity and, in particular,
\begin{equation}\label{eq:count}
 a_i-i=\#\{k:b_k<a_i\}.
\end{equation}

We prove~\eqref{eq:b-bound} by strong induction on $j$. The base case is
$b_1=1+2+3+4=10$. Fix $j\ge2$ and assume the bounds at every smaller index.
For $i\in\{j,j+1,j+2,j+3\}$, set $m=a_i-i$. Since $b_j>a_{j+3}\ge a_i$,
equation~\eqref{eq:count} gives $0\le m<j$.
If $m\ge1$, then $5m\le b_m<a_i$, which yields
\[
 a_i\le U_i=\left\lfloor\frac{5i-1}{4}\right\rfloor.
\]
For $m=0$, this inequality follows from $a_i=i$.

If $m+1<j$, then $a_i<b_{m+1}\le5(m+1)+5$; integrality gives
$4a_i\ge5i-9$, or
\[
 a_i\ge L_i:=\left\lfloor\frac{5i-6}{4}\right\rfloor.
\]
For the remaining case $m=j-1$, we have $a_i=i+j-1$ and
\[
 4a_i-(5i-9)=4j+5-i\ge3j+2>0,
\]
so the same lower estimate holds. Finally, for every integer $x$,
$\sum_{t=0}^3\lfloor(x+t)/4\rfloor=x$. Consequently,
\[
 \sum_{i=j}^{j+3}L_i=5j,
 \qquad \sum_{i=j}^{j+3}U_i=5j+5.
\]
Summing the estimates on $a_i$ proves the induction step.
\end{proof}

Define
\[
 d_j=b_j-5j\quad(j\ge1),\qquad d_0=b_0=0,
 \qquad e_n=U_n-a_n.
\]
Here index zero is only boundary notation. Lemma~\ref{lem:bounds} and strict
increase give
\begin{equation}\label{eq:d}
 0\le d_j\le5,\qquad d_{j+1}\ge d_j-1\quad(j\ge0).
\end{equation}
The inequality at $j=0$ follows directly from $d_1=5$.

\section{An exact error formula}

\begin{lemma}\label{lem:error}
The first four errors are zero. For $q\ge1$ and $1\le s\le4$,
\begin{equation}\label{eq:error}
 e_{4q+s}=\ind_{\{d_q\ge s\}}+
           \ind_{\{s=1,\ d_{q-1}=5\}}.
\end{equation}
Furthermore,
\begin{equation}\label{eq:window}
 d_j=5-(e_j+e_{j+1}+e_{j+2}+e_{j+3}).
\end{equation}
\end{lemma}

\begin{proof}
Write $u=U_{4q+s}=5q+s$. If $d_q<s$, then
$b_q<u<b_{q+1}$. Thus $u$ is selected and exactly $q$ forbidden integers are
below it; its selected index is $u-q=4q+s$. The error is zero. In this case
$d_{q-1}=5$ cannot coexist with $s=1$, by~\eqref{eq:d}.

Suppose $d_q\ge s$. For $s\ge2$,
\[
 b_{q-1}\le5q<u-1<b_q.
\]
The convention $b_0=0$ covers $q=1$. Therefore $u-1$ is selected with exactly
$q-1$ forbidden values below it; its selected index is again $4q+s$, giving
error one.

For $s=1$ and $d_q\ge1$, if $d_{q-1}<5$, then
$b_{q-1}<5q<b_q$, and the same count gives error one. If $d_{q-1}=5$, then
$q\ge2$, $b_{q-1}=5q$, and $d_q\ge4$. The bounds on the preceding forbidden
term, or $b_0=0$ for $q=2$, give
\[
 b_{q-2}<5q-1<b_{q-1}.
\]
Hence $5q-1$ has selected index $5q-1-(q-2)=4q+1$, giving error two.
This proves~\eqref{eq:error}. Since the four consecutive values
$U_j,\ldots,U_{j+3}$ sum to $5j+5$, their errors sum to $5-d_j$, proving
\eqref{eq:window}.
\end{proof}

\section{The two extreme sets}

Put
\[
 F=\{j\ge1:d_j=5\},\qquad Z=\{j\ge1:d_j=0\}.
\]
\begin{lemma}\label{lem:sets}
These sets satisfy
\begin{align}
 Z&=\bigcup_{k\in F}[4k+2,4k+5]_{\mathbb Z},\label{eq:Z}\\
 F&=\{1\}\cup\bigcup_{k\in Z}[4k-2,4k+1]_{\mathbb Z}.
 \label{eq:F}
\end{align}
\end{lemma}

\begin{proof}
Every four-term error window contains at most one error equal to two, because
such an error occurs only at an index congruent to one modulo four. Its sum
is five exactly when it contains a two and three ones.
By~\eqref{eq:error}, a two occurs at $4q+1$ exactly when $d_{q-1}=5$.
Then $d_q\ge4$, by~\eqref{eq:d}. The three preceding errors and the three
following errors are all one. Thus the four windows containing this two all
have sum five, and their starting indices are $4q-2,\ldots,4q+1$.
Equation~\eqref{eq:window}, with $k=q-1$, proves~\eqref{eq:Z} in both
directions.

For~\eqref{eq:F}, a value $d_j=5$ is equivalent to four consecutive zero
errors beginning at $j$. The initial zero window contributes $1\in F$.
Every four-term window starting at $j\ge2$ contains exactly one index
$4q+1$ with $q\ge1$. The error at this index is zero exactly when $d_q=0$.
In that case~\eqref{eq:d} implies $d_{q-1}\le1$. Formula~\eqref{eq:error}
then makes all three preceding and all three following errors zero. The
four windows containing the central index therefore start at
$4q-2,\ldots,4q+1$, and these are all the remaining positions in $F$.
\end{proof}

\section{The expanding recurrence and the runs}

Substitution of~\eqref{eq:Z} into~\eqref{eq:F} gives
\begin{equation}\label{eq:recurrence}
 F=\{1\}\cup\bigcup_{k\in F}[16k+6,16k+21]_{\mathbb Z}.
\end{equation}
The four consecutive indices $4k+2,\ldots,4k+5$ generate adjacent
four-integer intervals, whose union has precisely these endpoints.

Every generated index is larger than its parent. Thus any member of $F$
other than $1$ descends through smaller parents to $1$, and repeated
expansion of $\{1\}$ produces every member of $F$. Starting at
$I_0=[1,1]$, write $I_m=[\ell_m,r_m]$ for the interval at generation $m$.
Consecutive parents generate adjacent sixteen-integer intervals, so
\[
 \ell_{m+1}=16\ell_m+6,\qquad r_{m+1}=16r_m+21.
\]
Solving yields
\begin{equation}\label{eq:endpoints}
 \ell_m=\frac{7\cdot16^m-2}{5},\qquad
 r_m=\frac{12\cdot16^m-7}{5}.
\end{equation}
The endpoints are integers because $16^m\equiv1\pmod5$. Each interval has
length $16^m$, and
$\ell_{m+1}-r_m=20\cdot16^m+1>1$; consequently these are exactly the
maximal intervals in $F$.

\begin{proof}[Proof of Theorem~\ref{thm:runs}]
By~\eqref{eq:window}, a maximal zero-error run $[u,v]$ of length at least
four corresponds to the maximal interval $[u,v-3]$ in $F$. Conversely, a
maximal interval $[\ell,r]$ in $F$ gives zero errors on $[\ell,r+3]$; a
zero just outside either endpoint would extend the interval in $F$.
Applying this bijection to~\eqref{eq:endpoints} gives~\eqref{eq:runs},
including maximality and the assertion that all other runs have length
at most three.
\end{proof}

\section{Verification and research status}

An exact implementation of the original greedy definition was checked
through $400\,000$ selected terms. The error formula, both extreme-set
identities, and the expanding recurrence agreed throughout every available
complete window. Its only completed maximal equality runs of length at least
four were the five intervals displayed above. These finite checks support
implementation consistency; the argument above proves the infinite statement.

The sequence and the recorded conjecture are attributed to Hoang Xuan
Thanh. This note supplies a precise long-run statement and its proof, but
makes no claim of priority. The literature check is limited. The proof,
manuscript, and verification were developed with substantial assistance
from OpenAI's ChatGPT; additional AI-assisted mathematical checks were made
within the working session. This version has not received external human
peer review or proof-assistant certification.

\begin{thebibliography}{9}
\small
\bibitem{oeis}
H.~X. Thanh, \emph{A394676}, The On-Line Encyclopedia of Integer Sequences,
entry dated May 11, 2026.
\url{https://oeis.org/A394676}. Consulted November 27, 2026.
\bibitem{bounds}
M.~V. Prizmic, \emph{Uniform bounds for greedy sequences avoiding consecutive
sums}, research note, version 0.1, November 27, 2026.
\url{https://github.com/markprizmic1-source/greedy-consecutive-sums}.
AI-assisted manuscript; priority unverified.
\end{thebibliography}
\end{document}
